% IV. Програмна реализация.
% Пише се СЛЕД като кодът съществува. Засега: приетата структура и конвенции.
\section{Програмна реализация}
\label{sec:implementation}

Реализацията е на C++20 и се изгражда с CMake. Проектът е разделен на библиотека,
която съдържа цялата логика, и на изпълними файлове, които са тънки обвивки над нея.
Разделянето е направено заради тестовете: всичко, което заслужава проверка, живее в
библиотеката и се вика пряко от тестовете, без да се минава през команден ред.

\subsection{Обща структура на изходния текст}
\label{subsec:structure}

Публичните заглавни файлове стоят в \code{include/}, реализацията в \code{src/},
тестовете в \code{tests/}. Всеки от четирите слоя от Раздел~\ref{sec:design} е
отделен превод и отделна цел за компилация, така че зависимостите между слоевете да
са видими в конфигурацията, а не само в текста. Пълната йерархия на модулите и
списъкът на класовете се дават в \TODO{таблица с модулите, целите за компилация и
зависимостите между тях}.

\subsection{Разчитане на ръкостискането}
\label{subsec:impl_parser}

Разчитането е реализирано като поредица от функции, всяка от които приема изглед към
байтове и връща или разчетената структура, или причина за отказа. Няма изключения по
пътя на разчитането. Причината е, че повреденият или нарочно изкривен вход е нормално
явление за този слой, а не изключителна ситуация, и обработката му трябва да е
еднакво евтина при всеки пакет.

Всяка стъпка проверява дължината преди четенето. Това е основното изискване към слоя,
който работи с чужд вход, и е предмет на отделен набор от тестове с нарочно
подготвени неправилни съобщения. Резултатите от тези тестове се описват в
\TODO{резултати от проверката с неправилни и съкратени съобщения}.

\subsection{Проверка на веригата}
\label{subsec:impl_validator}

Изграждането на пътя е реализирано с обхождане в дълбочина по множеството от
кандидат-издатели, с ограничение върху дължината на пътя и с отчитане на вече
посетените, за да не се получи безкраен цикъл при кръстосано подписани сертификати.
Всяко отхвърляне на кандидат се записва с причина, така че при неуспех системата да
може да обясни защо нито един път не е бил приет, вместо просто да върне отказ.

Криптографските проверки на подписа се възлагат на външната библиотека. Всичко
останало, включително проверката на ограниченията върху имената, е собствена
реализация по \cite{rfc5280}. Конкретната библиотека и нейната версия се посочват в
\TODO{криптографска библиотека и версия, използвани в окончателната реализация}.

\subsection{Механизми за допускане}
\label{subsec:impl_admission}

Четирите механизма са реализирани зад общ интерфейс с едно решение: допускане или
отказ, придружено с причина. Веригата се конфигурира при стартиране и не се променя
по време на измерване, за да не се смесват режими в един и същ набор от данни.

Отчитането на таблицата с връзките използва структура, разделена по нишки, за да не
се превърне самото измерване в тясно място. Приетият подход и мотивите за него се
описват в \TODO{избраната структура за отчитане и мотивите за нея}. Избраните
листинги от изходния текст се дават в Приложение~\ref{sec:appendix}.
